\scanpage{154}
foncteur $M\mapsto M(-1)$ \uline{soit pleinement fidèle}.
Ses itérés sont les $M\mapsto M(-n)=M\otimes\Lambda(-n)$, $\Lambda(-n)\simeq\Lambda(\unsure{1|-1})^{\otimes n}$.
Noter que
\[
\operatorname{Hom}(\Lambda(-n),\Lambda(-n))\simeq\operatorname{Hom}(\Lambda(0),\Lambda(0))\simeq\mathbf Q
\]
On introduit alors formellement aussi les $M(n)$, $n\in\mathbf Z$, en considérant la catégorie pseudo-$\varinjlim$ des copies de $\mathcal M$, par les foncteurs $M(-1)$. Ainsi tout objet de la nouvelle catégorie $\widetilde{\mathcal M}$ \deleted{``motifs \unsure{virtuels}\inserted{\deleted{\unsure{effectifs}}}''} s'écrit
\[
M(i)\qquad i\in\mathbf Z,\quad M\in\operatorname{Ob}\mathcal M,\quad\text{avec}
\]
\[
\operatorname{Hom}(M(i),N(j))\simeq\operatorname{Hom}(M(i+r),N(j+r))
\]
où $r$ est pris tel que $i+r,j+r\leq0$, alors $M(i+r)$, $N(j+r)$ ont déjà un sens.
\margin{$\Lambda(1)$ de degré $-\rho$}
La catégorie $\widetilde{\mathcal M}$ est encore graduée (de $-\infty$ à $+\infty$ maintenant) périodique de période $\rho$ par le foncteur \deleted{\unsure{M}} $M\mapsto M(-1)$. En dim $i\geq0$ (\unsure{cas auquel} on se \unsure{ramène}) elle s'écrit comme
\[
\operatorname{Ps}\varinjlim_k\mathcal M^{i+k\rho},
\]
où
\diagram{\begin{tikzcd}
\mathcal M^{i+k\rho}\arrow[r] & \mathcal M^{i+(k+1)\rho}
\end{tikzcd}}
est le foncteur pl. fd $M\mapsto M(-1)$.
\margin{\uline{N.B.} Un sous motif d'un motif effectif est effectif..}
\mspage{153}
